problem:  
Let \( (M^n, g) \) be a complete, noncompact Riemannian manifold satisfying the following:  

1. The sectional curvature \( \mathrm{sec}_M \ge 0 \) outside a compact subset \( K \subset M \);  
2. For some basepoint \(p \in M\), all rescalings \( (M, r_i^{-2} g, p) \) with \(r_i \to \infty\) converge in the pointed Gromov–Hausdorff sense to a **unique tangent cone at infinity** \( C_\infty(M,p) = C(\Sigma) \);  
3. The link \( (\Sigma, g_\Sigma) \) of the cone is a smooth manifold with **\(\mathrm{Ric}_\Sigma \ge \lambda g_\Sigma > 0\)** (so it has strictly positive Ricci curvature in the induced cone metric).

**Question:**  
Does the asymptotic Ricci positivity of \( \Sigma \) imply that \( M \) has **finite topological type**?  
Equivalently, is \( M \) diffeomorphic outside a compact set to the **cylindrical end** \( \Sigma \times [0,\infty) \)?  

More strongly, if two complete, noncompact manifolds \( (M_1,g_1) , (M_2,g_2) \) satisfy these assumptions and share the same positively Ricci curved link \( (\Sigma, g_\Sigma) \) as their unique tangent cone at infinity, must their complements of compact subsets be **diffeomorphic**?  

Why is it a "good" problem:  
This formulation distills a core geometric–topological question that remains open even in simple settings: can strong **Ricci positivity at infinity** rigidify the topology of noncompact manifolds? It asks whether the **geometry of the cone link**—a purely asymptotic object—dictates not merely metric but also **topological** properties of the manifold at large scale.  

The problem unites three powerful frameworks:  
• **Cheeger–Colding theory** of metric measure convergence (controlling tangent cones at infinity),  
• **Ricci curvature geometry**, central to comparison and rigidity results,  
• **Riemannian topology of noncompact manifolds** (finite topological type, structure of ends).  

If resolved, the result would constitute an “asymptotic Ricci rigidity theorem,” showing that Ricci positivity at infinity suppresses topological complexity just as global nonnegative sectional curvature does in the classical Soul Theorem. The statement is simple to express—relating only \(M\), its asymptotic cone, and the geometry of \(\Sigma\)—but solving it would require deep new interactions between geometric analysis and topology. This blend of simplicity, plausibility, and potential foundational impact makes it a good research-level problem.